\begingroup\itshape
\begin{mdframed}[linewidth=1pt,linecolor=gray!60,backgroundcolor=gray!6,
  innertopmargin=12pt,innerbottommargin=12pt,innerleftmargin=14pt,innerrightmargin=14pt]
\section*{\textit{Resumen ejecutivo}}
El Centro Javeriano de Competitividad presenta diez propuestas sobre productividad, inteligencia artificial (IA) y financiación empresarial para el Plan Nacional de Desarrollo 2026–2030.

\begin{enumerate}[label=\arabic*.,leftmargin=*,itemsep=4pt,parsep=0pt,topsep=6pt]
\item \textbf{Adopción de IA en mipymes.} Acompañar a micro, pequeñas y medianas empresas en el diagnóstico de procesos, la selección y la implementación de herramientas.
\item \textbf{Investigación con participación empresarial.} Financiar alianzas de empresas y universidades con dinero, horas de personal, materiales o uso de equipos, registrando los recursos que cada empresa entrega.
\item \textbf{Educación y reconversión laboral.} Formar a estudiantes y docentes en IA y ofrecer bonos para cursos laborales que evalúen lo aprendido en tareas de cada ocupación.
\item \textbf{Cofinanciación condicionada de IA a empresas usuarias.} Cofinanciar proyectos conjuntos y vincular los desembolsos a entregas de cada etapa, como datos preparados, pruebas técnicas o herramientas en operación.
\item \textbf{Servicios profesionales y administrativos.} Ayudar a firmas contables, jurídicas y administrativas a usar IA con registro de fuentes, revisión profesional y protección de datos de clientes.
\item \textbf{Pruebas financieras supervisadas.} Probar aplicaciones de IA bajo supervisión de la Superintendencia Financiera (SFC), con límites de duración y usuarios y condiciones para suspenderlas.
\item \textbf{Acceso a cómputo.} Cofinanciar acceso temporal a equipos compartidos o servicios de procesamiento de datos que el proyecto necesite para funcionar.
\item \textbf{Capital y diversificación del riesgo.} Eliminar la garantía de rentabilidad mínima pensional, evaluar el riesgo y el rendimiento de la cartera en su conjunto y fijar en \$2.000 millones el mínimo de capital para administrar pensiones.
\item \textbf{Herramientas compartidas en el Estado.} Comprar conjuntamente herramientas que varias entidades puedan usar para buscar documentos, clasificar solicitudes y preparar borradores.
\item \textbf{Proyectos con resultados verificables.} Seleccionar proyectos con presupuesto y responsables y comprobar si reducen tiempos, errores o costos o mejoran la calidad del servicio.
\end{enumerate}
\end{mdframed}
\par\endgroup
